\printedpage{140}
% head: A. Grothendieck : Dérivateurs, Ch. XII, Édition des Universités de Montpellier II et Paris VII

\section*{5 Axiomatics of non-locally constant relative $W$-homotopy types $\mathrm{Ho}^{\mathcal{L}}_W(S)$.}

We have said nothing up to this point about possible relative $W$-homotopy types (over an object $S$ of $\mathcal{M}$) which would \textit{not} be ``locally constant''. Yet they were discussed in loc. cit. (VI, theorem-scholie 2, part II of the theorem pages 109-113.) I take the opportunity to come back to this. We must introduce, in each $\mathcal{M}/S$, a second localizer
\[
W^0_S \subseteq \mathrm{Fl}(\mathcal{M}/S), \qquad W^0_S \subseteq W_S
\]
(finer than $W_S$), which will make it possible to define
\msp{page 219}
\[
\mathrm{Ho}_W(S) \deff (\mathcal{M}/S)(W^0_S)^{-1}.
\]
That there is no \textit{canonical} choice which imposes itself here, deduced from the structure $(\mathcal{M},W)$ alone, is clear from the example of $(\mathrm{Cat},W)$, in which case we find \textit{two} symmetric choices, transformed into one another by the involution
\[
X \longmapsto X^\circ : (\mathrm{Cat},W) \xrightarrow{\sim} (\mathrm{Cat},W),
\]
namely $W^{\mathrm{g}}_S$ and $W^{\mathrm{d}}_S$. The choice of one or the other, say (to fix ideas) of $W^{\mathrm{g}}_S$, is associated with a notion of \textit{local isomorphism} (or of ``local co-isomorphism'' in the case of $W^{\mathrm{d}}_S$) which I would like to axiomatize. It does seem that for the main categories of models for Hot (and for its variants), we have in a fairly natural way such a notion of local isomorphism (without the encumbrance, most often, of an involutive automorphism providing a symmetric notion \dots).

We shall therefore suppose, as in loc. cit. p. 94-96, given a class
\[
\mathcal{L} \subseteq \mathrm{Fl}(\mathcal{M})
\]
of morphisms, playing the role of ``arrows of % [sic] jounant -> jouant, faute de frappe
\msp{page 220}
localization'' or ``of local isomorphisms'', satisfying the following conditions.

\begin{enumerate}
\item[$(*)$] $\mathcal{L}$ contains the isomorphisms, is stable under composition, and under base change.
\end{enumerate}

We put, for $S$ in $\mathcal{M}$,
\[
\mathcal{L}_S \subseteq \mathcal{M}/S \qquad \text{full subcategory of the } S'/S \text{ such that } S' \longrightarrow S \text{ \texttt{[is]} in } \mathcal{L}.
\]
We then define for every $S$ in $\mathcal{M}$
\[
\begin{aligned}
&W^{\mathcal{L}}_S \subseteq \mathrm{Fl}(\mathcal{M}/S)\\
&(X \xrightarrow{u} Y) \in W^{\mathcal{L}}_S \iff \forall\, S' \in \mathcal{L}_S,\ (u_{S'} : X \times_S S' \longrightarrow Y \times_S S') \text{ is in } W.
\end{aligned}
\]
% foot: Version 30 mars 2003

\printedpage{141}
% head: A. Grothendieck : Dérivateurs, Ch. XII, Édition des Universités de Montpellier II et Paris VII

This is therefore a variant of the localizer $W^{\mathrm{u}}_S$, where now we take as base changes only the ``$\mathcal{L}$-localizing'' arrows, i.e. $\in \mathcal{L}$. We thus have
\[
W^{\mathrm{u}}_S \subseteq W^{\mathcal{L}}_S \subseteq W_S
\]
as subsets of $\mathrm{Fl}(\mathcal{M}/S)$, and likewise in any full subcategory $\mathcal{F}$ of $\mathcal{M}/S$. In the case of $\underline{\mathrm{Fib}}_W S$, by virtue of axiom \circled{b} above ($W^{\mathrm{u}}_S = W_S$ in $\underline{\mathrm{Fib}}_W S$), the three localizers coincide.

We put
\[
\mathrm{Ho}^{\mathcal{L}}_W(S) \text{ (or } \mathrm{Ho}_W(S)\text{)} \ \deff (\mathcal{M}/S)(W^{\mathcal{L}}_S)^{-1},
\]
and we gave in loc. cit. (part II of the theorem-scholie) an interesting result on this invariant, for fixed $S$. My aim
\msp{page 221}
at present is to see what happens for variable $S$. Visibly, a notion of relative $W$-homotopy type over an arbitrary object $S$ of $\mathcal{M}$ is satisfactory, only if it allows inverse images by \textit{arbitrary} arrows
\[
f : S' \longrightarrow S
\]
in $\mathcal{M}$, which absolutely must define
\[
f^{*\mathcal{L}}_W \text{ or } f^*_W : \mathrm{Ho}^{\mathcal{L}}_W(S) \longrightarrow \mathrm{Ho}^{\mathcal{L}}_W(S').
\]
We are therefore faced with the same perplexity as just now - and it is \textit{not,} this time, the introduction of $\underline{\mathrm{Fib}}_W(S)$ that is going to save the situation, for
\[
(\underline{\mathrm{Fib}}_W S)(\underbrace{W^{\mathcal{L}}_S}_{\substack{= W_S \\ = W^{\mathrm{u}}_S}})^{-1} \longrightarrow (\mathcal{M}/S)(W^{\mathcal{L}}_S)^{-1}
\]
has no desire, this time, to be an equivalence of categories, but much rather an inclusion (fully faithful) of the relative $W$-homotopy types constant over $S$, into arbitrary $W$-homotopy types. An axiom is visibly needed to replace the ``homotopy fiber'' axiom.

\msp{page 222}

In the case of Cat, the class $\mathcal{L}$ of local isomorphisms (say) allows us to define, as above, the localizer $W^{\mathcal{L}}_S = W^{\mathrm{g}}_S$ in the $\mathrm{Cat}/S$, as consisting of the $u : X \longrightarrow Y$ over $S$ which are in $W$ and remain so under every base change $S' \longrightarrow S$ which is in $\mathcal{L}$. But there is reason to \textit{saturate} $\mathcal{L}$, by considering the class $\mathcal{L}^\wedge$ of arrows $f : S \longrightarrow S'$ of Cat such that the base change
\[
f^* : \mathrm{Cat}/S \longrightarrow \mathrm{Cat}/S'
\]
sends $W^{\mathrm{g}}_S$ into $W^{\mathrm{g}}_{S'}$. This $\mathcal{L}^\wedge$ is tautologically stable under composition, and contains the isomorphisms (better, contains $\mathcal{L}$ = set of local isomorphisms). It is not
% foot: Version 30 mars 2003
\printedpage{142}
% head: A. Grothendieck : Dérivateurs, Ch. XII, Édition des Universités de Montpellier II et Paris VII
clear, from the definition given here, that it is stable under base change - we are going to strengthen the definition accordingly, by requiring, for $f : S' \longrightarrow S$ to be in $\mathcal{L}^\wedge$, that for every $T$ over $S$, let us consider $f_T : T' = S' \times_S T \longrightarrow T$, $f^*_T : \mathrm{Cat}/T \longrightarrow \mathrm{Cat}/T'$ be compatible with the localizers $W^{\mathrm{g}}_T$, $W^{\mathrm{g}}_{T'}$. Then $\mathcal{L}^\wedge$ is (just like $\mathcal{L}$) stable under composition and under base change, and moreover it contains $\mathcal{L}$. Moreover, for every $S$, we have $W^{\mathcal{L}}_S = W^{\mathcal{L}^\wedge}_S$. This little sorites could of course be carried out for any modelizer $(\mathcal{M},W)$. But the interesting thing here is that $\mathcal{L}^\wedge$ contains the set $\mathrm{Liss}_W$ of $W$-smooth arrows. (I do not know whether it is equal.) We could therefore have
\msp{page 223}
just as well started from the set of arrows $\mathrm{Liss}_W$, to define the localizers $W^{\mathrm{g}}_S$ as
\[
W^{\mathrm{g}}_S = W^{\mathrm{Liss}_W}_S \qquad \text{instead of } W^{\mathrm{g}}_S = W^{\mathrm{Isoloc}}_S,
\]
which is equivalent.

On the other hand, in working with the localized categories $\mathrm{HOT}^\circ_W(S) = (\mathrm{Cat}/S)(W^{\mathrm{g}}_S)^{-1}$, the role of $\underline{\mathrm{Fib}}_W S$ in the theory of the $\mathrm{HOT}^{\mathrm{lc}}_W(S)$ is played here by $\underline{\mathrm{Prop}}_W(S)$. We thus see that in the theory of non-locally constant relative $W$-homotopy types, $\mathrm{HOT}^\circ_W$ mode, the sets of arrows $\mathrm{Liss}_W$ and $\mathrm{Prop}_W$ play a role that is in a sense complementary: the arrows $f : S' \longrightarrow S$ of $\mathrm{Liss}_W$ are (practically) those which allow us directly to define
\[
f^{*\circ}_W : \mathrm{HOT}^\circ_W(S) \longrightarrow \mathrm{HOT}^\circ_W(S')
\]
as induced by
\[
f^* : (\mathrm{Cat}/S, W^{\mathrm{g}}_S) \longrightarrow (\mathrm{Cat}/S', W^{\mathrm{g}}_{S'}),
\]
whereas the arrows of $\mathrm{Liss}_W$ allow us to define for every $S$ in Cat the full subcategory
\[
\underline{\mathrm{Prop}}_W S \subseteq \mathrm{Cat}/S,
\]
giving rise to an equivalence of categories
\[
\left\{
\begin{aligned}
&(\underline{\mathrm{Prop}}_W S)(W^{\mathrm{u}}_S)^{-1} \xrightarrow{\sim} \underbrace{(\mathrm{Cat}/S)(W^{\mathrm{g}}_S)^{-1}}_{= \mathrm{HOT}^0_W(S)} % ?? uncertain: exposant de HOT sous l'accolade -- 0 (zéro) ou \circ
\\
&\text{with}\\
&W^{\mathrm{u}}_S = W^{\mathrm{g}}_S \text{ on } \underline{\mathrm{Prop}}_W S
\end{aligned}
\right.
\]
(which visibly is the announced substitute for ``the $W$-homotopy fiber axiom''). This
\msp{page 224}
equivalence, which allows us to describe $\mathrm{Hot}^\circ_W(S)$ in terms of $\underline{\mathrm{Prop}}_W S \subseteq \mathrm{Cat}/S$ and a localizer $W^{\mathrm{u}}_S$ \textit{stable under} \textbf{every} \textit{base change,} allows us on the other hand to define the expected contravariant law for $\mathrm{HOT}^\circ_W(S)$.
% foot: Version 30 mars 2003

\printedpage{143}
% head: A. Grothendieck : Dérivateurs, Ch. XII, Édition des Universités de Montpellier II et Paris VII

On the other hand, when we interchange the role of $\mathrm{Liss}_W$ and $\mathrm{Prop}_W$, we find the complementary theory of $\mathrm{HOT}_W(S)$,
\[
\left\{
\begin{aligned}
&\mathrm{HOT}_W(S) \deff (\mathrm{Cat}/S)(W^{\mathrm{d}}_S)^{-1} \xleftarrow{\sim} (\underline{\mathrm{Liss}}_W S)(W^{\mathrm{u}}_S)^{-1},\\
&\text{on } \underline{\mathrm{Liss}}_W S,\ W^{\mathrm{d}}_S = W^{\mathrm{u}}_S,
\end{aligned}
\right.
\]
the description of $\mathrm{HOT}_W(S)$ in terms of $\underline{\mathrm{Liss}}_W S$ and $W^{\mathrm{u}}_S$ allows us to define its contravariant character in $S$, for \textit{every} arrow $f : S' \longrightarrow S$ in Cat. Whereas if we want to be able to pass to the categories of fractions on $\mathrm{Cat}/S$ and $\mathrm{Cat}/S'$ in their entirety (via $W^{\mathrm{d}}_S$, $W^{\mathrm{d}}_{S'}$), we must restrict ourselves to $W$-proper base changes.

So the perplexity comes to me whether this complementary, or even this duality, between two sets of arrows, the ``smooth'' and the ``proper'', might not be a general fact which turned up at least in the categories of models close to the description of HOT or of a $\mathrm{Hot}_W$.

\msp{page 225}

It may well be that in $\boldsymbol{\Delta}^\wedge$ we have a perfect duality between these two kinds of notions (I shall have to examine the matter), perhaps even in any category $\mathcal{A}^\wedge$ (but in $\boldsymbol{\Delta}^\wedge$ there is an involutive automorphism $K \longmapsto K^\circ$, corresponding to the $X \longmapsto X^\circ$ in Cat \dots). And after all, even in the context of topological spaces, the notions of properness on the one hand, of local asphericity on the other, for an arrow $X \longrightarrow Y$, are well defined, and might well fit together as do the corresponding notions in Cat, and (who knows) give rise to two complementary versions, or ``dual'' to one another, of a notion of relative homotopy type.

Not being clear about this, I am a little undecided how to present an axiomatic description of the theory of the $\mathrm{Ho}^{\mathcal{L}}_W(S)$, for a general category of models $(\mathcal{M},W)$. Should I give myself from the outset
\[
\mathrm{Liss}, \mathrm{Prop} \subseteq \mathrm{Fl}(\mathcal{M}),
\]
stable under composition, under base change, containing the isomorphisms, and such that (perhaps)
\[
\mathrm{Liss} \cap \mathrm{Prop} \subseteq \mathrm{Fib}_W,
\]
\msp{page 226}
but above all they allow us to define, for $S$ in $\mathcal{M}$
\[
\left\{
\begin{array}{lcccl}
W^{\mathrm{g}}_S & \deff & W^{\mathrm{Liss}}_S & \subseteq & \mathcal{F}(\mathcal{M}/S)\\
W^{\mathrm{d}}_S & \deff & W^{\mathrm{Prop}}_S & \subseteq & \mathcal{F}(\mathcal{M}/S)
\end{array}
\right.
\]
with the following properties:

\begin{enumerate}
\item[{\circled{a}}] If $f : S' \longrightarrow S$ is in Prop (resp. Liss), then $f^* : \mathcal{M}/S \longrightarrow \mathcal{M}/S'$ sends $W^{\mathrm{d}}_S$ into $W^{\mathrm{d}}_{S'}$ (resp. $W^{\mathrm{g}}_S$ into $W^{\mathrm{g}}_{S'}$),
\end{enumerate}
% foot: Version 30 mars 2003
\printedpage{144}
% head: A. Grothendieck : Dérivateurs, Ch. XII, Édition des Universités de Montpellier II et Paris VII
\noindent which is a pure tautology, and does not yet concern the relations of Prop with Liss, but on the other hand:

\begin{enumerate}
\item[{\circled{b}}]
\[
\left\{
\begin{aligned}
&(\underline{\mathrm{Liss}}S)(W^{\mathrm{u}}_S)^{-1} \xrightarrow{\sim} (\mathcal{M}/S)(W^{\mathrm{d}}_S)^{-1} \text{ \texttt{[is]} an equivalence}\\
&\text{and } W^{\mathrm{u}}_S = W^{\mathrm{d}}_S \text{ on } \underline{\mathrm{Liss}}S,
\end{aligned}
\right.
\]
and dually
\[
\left\{
\begin{aligned}
&(\underline{\mathrm{Prop}}S)(W^{\mathrm{u}}_S)^{-1} \xrightarrow{\sim} (\mathcal{M}/S)(W^{\mathrm{g}}_S)^{-1} \text{ \texttt{[is]} an equivalence}\\
&\text{and } W^{\mathrm{u}}_S = W^{\mathrm{g}}_S \text{ on } \underline{\mathrm{Prop}}S.
\end{aligned}
\right.
\]
\end{enumerate}

Moreover, there is reason to state \circled{c} two adjunction properties (dual to one another), for which I refer to VI, theorem-scholie 2, p. 108-113. Here the subsidiary question arises of simplifying the conditions of validity of the adjunction formula, somewhat disparate in loc. cit.

On the other hand, the question arises of developing the theory of the covariant laws $f_!$ and $f_*$ for $\mathrm{Ho}^{\mathrm{lc}}_W(S)$ and for $\mathrm{Ho}^{\mathrm{Prop}}_W(S)$, $\mathrm{Ho}^{\mathrm{Liss}}_W(S)$.

\msp{page 227}

I have good hope that the functor $f_!$ can be described, and its adjunction with $f^*$ established, at as little cost in this general context as in Cat. But I do not feel like dwelling on it at present.

\section*{6 $W$-homotopy cofiber in $\mathrm{Cat}$.}

On the other hand, I feel like examining here under what conditions on the basic localizer $W$, there exists in $(\mathrm{Cat},W)$ a theory dual to that of the $W$-homotopy fiber.

The key question is, of course, that of having a good grip on $W$-cofibrations. Now so far I know \textit{only} the open Dwyer immersions, the closed Dwyer immersions, and what follows from them by composition etc. The work of \textsc{Thomason} may motivate the hope that knowledge of those immersions, which seem here to play a role similar (dual) to that of $W$-perfect morphisms, will suffice to develop a good theory of ``$W$-homotopy cofibers''.

% trois traits verticaux imprimés dans la marge gauche de ce paragraphe
{\reversemarginpar\marginnote{\textbar\textbar\textbar}}I assume throughout this section that $W$ satisfies W(7 bis) (axiom of cocartesian squares), to have good properties of open Dwyer immersions (\footnote[124]{\begin{enumerate}
\item[a)] These are $W$-cofibrations (follows from b) and c)).
\item[b)] If
\end{enumerate}}), and more generally, for $W$-cofibrations. % ?? uncertain: coupure en fin de ligne "co-cartésiens" -- cocartésiens ou co-cartésiens
% foot: Version 30 mars 2003